\chapter{Discussion}

\section{Communautés phytoplanctoniques de la région d'étude}
\subsection{Distributions des contributions des pigments photo-synthétiques}
Les résultats présentés par la figure \ref{fig:violininterannualratiospigmenttotpig} ne montrent pas de tendances claire et significative d'augmentation ou de diminution de la contribution d'un groupe phytoplanctonique en particulier. Ces observations vont dans le sens de celles émises par Adroli dans sa publication récente \cite{adroli2026} : le secteur indien semble moins impacté par des changements de structuration des communautés phytoplanctoniques que les autres bassins de l'océan austral. Ces résultats sont cependant à mettre en perspectives avec le nombre de données effectivements utilisées pour calculer les tendances, variable d'une année à une autre (voir annexe \ref{annexe:n_data_distribs}). De plus, la couverture nuageuse importante de la région ne permet de recueillir en moyenne que 10 à 20\% de données par satellite, des biais d'échantillonages peuvent donc être présents. D'autre part, des changements de structuration des communautés phytoplanctoniques peuvent avoir lieu à plus fine échelle et ne pas être détectés par ces calculs de distributions. Des tendances interanuelles ont été effectuées par fod (voir annexe \ref{fig:annexe_violins_all}), aucune tendance forte n'a été détécté. Enfin, pour réellement évaluer des changements de structuration de communautés phytoplanctonique, un jeu de donnée réparti sur au moins 40 ans est nécessaire \cite{thomalla2023Widespread}.


\subsection{Structure des communautés phytoplanctoniques}
Les clusters de communautés phytoplanctoniques présentés par les figures \ref{fig:distributionpigmentscommunauteskmeansk6pftstats} et  \ref{fig:distribpigmentsincommunauteskmeansk6pftstats} présentent des groupes de données ayant des contributions de pigments et une biomasse totale similaire (la variable de chlorophylle-a totale est incluse dans la PCA et le clusstering). L'équivalence directe entre les ratios de concentration de pigments photo-synthétiques et les types/classe de taille du phytoplancton est seulement hypothétique et doit être interprétée avec précaution. Des mesures \textit{in-situ} par HPLC et Cytométrie en flux sont nécessaires afin de qualifier exactement la composition taxonomique des communautés phytoplanctoniques. Cette simplification reste néanmoins informative. Les clusters projetés sur une carte présentent une structuration spatiale (figure \ref{fig:mappuritycommunauteskmeansk6pft}). Cette observation appuie l'hypothèse que, malgré les simplifications et la réduction de dimension, la méthode d'évaluation des communautés phytoplanctonique adoptée permet une représentation pas trop éloignée de la réalité biologique documentée. Les assemblages de PFT associés à chaque cluster sont cohérents avec la littérature : la communauté 4, dominée par les pigments dvChla et Zeaxanthine, est présent uniquement au Nord de la région, ce qui est cohérent avec les travaux récents qui montrent l'absence de Prochlorococcus (Cyanobactérie) à des latitudes très faibles \cite{Pan2020}, \cite{Johnson2006}. La forte biomasse observée autour de Kerguelen est attendue : la concentration en fer du milieu, normalement limitée dans la région, augmente avec les apports continentaux et induit des blooms importants \cite{Blain2007}.
Il serait intéressant d'effectuer ces mêmes analyses sur une période plus longue afin d'évaluer les restructurations de ces communauté avec une plus grande certitude. Une analyse distincte de la biomasse totale et des contributions phytoplanctoniques serait également informative.
La biomasse de NTI n'est pas significativement différente selon la communauté phytoplanctonique associée. Cela peut s'expliquer par la phénologie connue : il existe un délai temporel entre bloom phytoplanctonique et colonisation par les NTI \cite{nicholson2025satellite}, \cite{kauko2021phenology}, \cite{lampe2026differences}. D'autre part, la quantité limitée de donnée dû a un échantillonage couteux, associé à la faible couverture spatio-temporelle limite la fiabilité de ces résultats. D'avantage de données et un échantillonage couvrant d'avantage la région sont nécessaire pour confirmer ces résultats. Par ailleurs, des analyses de distributions de NASC inter-annuelles et sur toutes la ROI ont été réalisés (voir Annexe \ref{annexe:distrib_nasc}) et n'ont pas montré de changements significatifs de biomasse de NTI à cette échelle.


\section{Structure prédictive du NASC}
Les résultats de prédiction des différents modèles, présentés dans la section \ref{res:trees_perf} présentent des performances prédictives contrastées entre les différents schémas de validation et fréquences. Les RMSE ne sont pas directement comparables entre fréquences puisque le NASC n'a pas la même plage de valeurs pour chacune d'entre elle. Une statistique comme la NRMSE (RMSE / écart-type du NASC observé la fréquence considérée) permettrait de comparer directement les prédictions aux deux fréquences. On remarque cependant visuellement une performance inférieure à 120 kHz qui pourrait s'expliquer comme suis : à 120 kHz, les données sont échantillonnées jusqu'à 300m et, en journée (période d'étude), les NTI sont majoritairement présents à des profondeurs supérieures à cause de la migration nycthémérale (DVM). Une grande partie du signal biologique est donc absente et le modèle n'arrive pas à apprendre la structure des données.

\subsection{Plafond d'apprentissage}
La figure \ref{fig:learning_curve} montre un plafond structurel. L'écart-type entourant la courbe est néanmoins important et indique une variabilité inter-fold conséquente. Les figures \ref{label}, présentées en annexe \ref{label}, présentent les courbes d'apprentissage par fold, pour chaque schéma de validation croisée. Pour chaque fold, le même motif se répète : la courbe train (bleue) est plate dès 10\% du train, indiquant que le modèle atteint son plafond de capacité quasi instantanément, indépendamment du volume de données disponible. On observe également une hétérogénéité marquée entre blocs : certains folds sont "faciles" : les données de test ressemblent suffisamment à l'ensemble d'entraînement pour que le modèle généralise très bien dès la plus petite fraction testée. D'autres folds sont "difficiles" : l'erreur de test reste élevée et aucune donnée supplémentaire ne permet au modèle de s'améliorer.

Il existe donc un facteur limitant important, dont deux hypothèses peuvent expliquer l'origine : (1) l'information contenue dans les covariables est insuffisante pour décrire la variabilité du NASC, (2) la difficulté d'extrapolation imposée par certains blocs est trop importante pour être surmontée avec les covariables actuelles. Si cette seconde hypothèse était effectivement vrai, le schéma Monte-Carlo (MC-CV), qui n'impose aucun blocage spatio-temporel, devrait présenter une courbe d'apprentissage décroissante, l'ajout de données réduisant progressivement la difficulté d'extrapolation résiduelle. Or ce n'est pas ce qui est observé (figure~\ref{label}) : la courbe MC-CV plafonne elle aussi, dès les plus petites fractions de train testées.

Une hypothèse alternative permet d'expliquer ce résultat pour le schéma naïf : les données présentent une forte redondance spatio-temporelle (une mesure de NASC toutes les 3 à 4 secondes), telle que, même à 20\% du train, l'échantillon aléatoire reste suffisamment dense pour qu'un quasi-voisin de chaque point de test soit déjà présent dans ce sous-échantillon. Le plafond observé refléterait alors une saturation par redondance plutôt qu'un véritable plafond de capacité du modèle.

Une analyse complémentaire, mettant en relation la RMSE de test avec la distance géographique moyenne au jeu d'entraînement pour chaque fold du schéma de blocage 20 km x 20km, montre une corrélation positive mais faible et non significative (annexe \ref{annexe:corr_dist}) : plusieurs folds à distance comparable présentent des performances très contrastées, et la taille du jeu d'entraînement, quasi constante entre folds, n'explique pas non plus cette hétérogénéité. Il est donc difficile, en l'état, de mettre en évidence la cause exacte du plafonnement observé : distance géographique sur l'axe latitudinal (non-stationnarité d'ordre 1), extrapolation dans l'espace des covariables, ou hétérogénéité de la structure spatiale locale (non-stationnarité d'ordre 2). Des analyses complémentaires, notamment en lien avec le respect des hypothèses de stationnarité d'ordre 1 et 2, sont nécessaires pour trancher entre ces hypothèses (voir section \ref{annexe:lat_effect}).

\subsection{Limites de chaque schéma de validation}
\subsubsection{MC-CV 80/20}
Le schéma de validation MC-CV a des performances de prédiction élevées. Ces résultats sont cependant à considérer avec précautions : les données environnementales étant proche spatio-temporellement, l'ensemble de validation n'est pas indépendant de l'ensemble d'entraînement (fuite d'information). Les performances prédictives évaluées par la RMSE/$R^{2}$ sont donc artificiellement optimistes et ne reflètent pas une vraie capacité prédictive du modèle. Au delà du simple risque de surestimation des performances prédictives, ce modèle présente possiblement de moins bonnes performances de généralisation en situation d'extrapolation que les autres : les hyperparamètres du modèle naive ont été sélectionnés par un tuning qui n'a jamais pénalisé la dépendance à un voisin proche, puisque cette dépendance était justement récompensée par la fuite. Le modèle naive n'a donc aucune raison d'avoir appris à extrapoler (contrairement au tuning fait en 1500×1000km qui sélectionne des hyperparamètres devant fonctionner correctement sans point proche disponible). 


\subsubsection{Résolution 1500 km x 1000 km}
Concernant la résolution 1500 x 1000 km, les folds sont peu nombreux (5) mais contiennent beaucoup de données ont une variance importante. Le modèle est supposé mieux généraliser dans ce cas, puisque son évaluation porte sur un contexte proche d'une véritable extrapolation à grande échelle. En contrepartie, le faible nombre de folds réduit la puissance statistique de l'estimation de performance, et chaque estimation individuelle de RMSE/R\textsuperscript{2} repose sur peu d'observations indépendantes. La figure \ref{label} évalue effectivement faiblement les performances de prédictions par comparaison avec les autres résolutions. En annexe \ref{}, une analyse plus fine par fold est présentée. On remarque que le modèle est fortement biaisé pour deux folds en particulier : s\_0\_-2 et s\_1\-2. En regardant la localisation spatiale du test et la distance géographique minimale entre train et test, on comprend que les données test sont potentiellement très éloignées de l'ensemble d’entraînement et donc que l'on demande excessivement au modèle d'extrapoler. En regardant les distances euclidiennes moyennes entre les covariables numériques du train et du test (toutes les covariables sauf FOD), on s'aperçoit que la distance n'est pas particulièrement élevée. En tout cas, pas plus que pour d'autres folds où les performances prédictives sont supérieures. Ce n'est qu'en regardant la moyenne de chaque covariable par ensemble de donnée (train/test) et en mettant en relation avec l'importance des variables évaluée section\ref{label} que l'on s'aperçoit des écarts importants entre des variables clés et de la demande d'extrapolation du modèle. Pour le fold s\_0\_-2, une forte biomasse de phytoplancton est observée (chlorophylle-a totale importante), associée à une proportion de diatomées et de pélagophytes importante (resp. fucoxanthine et 19'butanoyloxyfucoxanthine élevée) et une proportion d'haptophyte remarquablement faible (19'-hexanoyloxyfucoxanthine). Le modèle sous-estime le NASC, ce qui est cohérent avec l'analyse des covariables effectuée : la biomasse de phytoplancton est particulièrement importante localement et la biomasse de NTI croit en conséquence. 

\subsubsection{Résolution 200 km x 200 km}
Cette résolution intermédiaire permet d'isoler 24 blocs éligibles pour RF, conduisant respectivement à 7 et 23 folds retenus. Le nombre de folds, plus élevé qu'à 1500 x 1000 km, offre une meilleure puissance statistique pour estimer la performance moyenne, tout en imposant une extrapolation spatiale à une échelle encore relativement sévère.
Les performances observées (RMSE test moyen d'environ 0,45 pour Random Forest à 38 kHz) restent proches de celles obtenues à 1500x1000 km, suggérant que la difficulté d'extrapolation demeure élevée à cette échelle intermédiaire. Le nombre encore restreint de folds à cette résolution implique qu'une part de la variabilité observée entre folds peut provenir de la taille d'échantillon de test réduite par fold plutôt que d'une réelle hétérogénéité de difficulté spatiale.

\subsubsection{Résolution 20 km x 20 km}
À cette résolution fine, 106 blocs sont éligibles pour RF, ce qui conduit à un plafonnement à 30 folds (borne maximale). Le buffer associé (5 km) étant faible, les blocs testés restent géographiquement proches du reste du jeu d'entraînement conservé, la tâche imposée au modèle se rapproche donc davantage d'une interpolation que d'une véritable extrapolation à grande échelle. Le nombre de données par fold est cependant restreint et la variance diminue. De manière attendue, les performances obtenues (RMSE test moyen d'environ 0,27 pour Random Forest à 38 kHz) sont sensiblement meilleures qu'aux deux résolutions plus strictes. L'analyse des distributions de valeurs prédites à cette résolution (cf. figure \ref{fig:distrib_nasc_38_reel_pred}) révèle cependant une concentration marquée des prédictions autour de quelques valeurs modales. L'homogénéité des covariables à l'intérieur des folds limite donc la capacité du modèle à différencier finement les observations d'un même fold, indépendamment de son réglage. 

\subsubsection{Résolution journalière}
Le blocage temporel isole 19 blocs éligibles pour RF, conduisant respectivement à 6 folds retenus. Contrairement aux résolutions spatiales précédentes, ce schéma évalue la capacité du modèle à généraliser vers une journée entière dont aucune observation n'a été vue à l'entraînement, plutôt que vers une région géographique. 
Les performances obtenues (RMSE test moyen d'environ 0,43 pour Random Forest à 38 kHz) se situent à un niveau inférieure de toutes les résolutions spatiales considérées, suggérant que l'extrapolation vers une journée non observée constitue une difficulté comparable à une extrapolation spatiale à échelle, même grande.
La calibration évaluée en figure \ref{label} à 120kHz montre un fort bruit et une incapacité du modèle apprendre les structures discriminantes. Ces résultats sont néanmoins à mettre en perspective avec le faible nombre de données par quartile pouvant bruiter les résultats.  


\subsection{Effet de la latitude} 
\label{annexe:lat_effect}
La structuration des données selon le gradient latitudinal n'a pas été directement considérée lors de la construction du modèle. Seuls les FOD capturent une partie de ce gradient. Le non-respect de la stationnarité d'ordre 1, combiné à l'absence d'une covariable capturant pleinement ce gradient (par exemple une covariable de latitude explicite), peut affecter l'importance des variables : une covariable fortement corrélée à la latitude peut apparaître comme très importante non pas parce qu'elle apporte une information écologique directe, mais parce qu'elle encode indirectement la position latitudinale, ce qui pourrait être le cas ici pour le FOD. Une amélioration possible du modèle serait de détendancer les données en amont (régression de la variable réponse sur la latitude, puis modélisation des résidus de cette régression)

Le non-respect de la stationnarité d'ordre 2 induit des conséquences moins visibles : la structure même de la corrélation spatiale (portée et forme de l'autocorrélation) peut elle-même varier selon la région considérée. Par exemple, une portée de corrélation plus courte à proximité des îles Kerguelen, où les gradients océanographiques sont resserrés, que dans les zones de plein océan. Une unique résolution et buffer, appliqué uniformément à toute la zone d'étude quelle que soit la résolution, n'offre alors pas la même protection contre la fuite d'information partout : trop large dans les régions à faible portée réelle d'autocorrélation, ce qui réduit inutilement le jeu d'entraînement disponible ; trop court dans les régions à forte portée, ce qui laisse subsister une fuite résiduelle malgré le buffer. Cette hétérogénéité pourrait expliquer, au moins en partie, la forte variance inter-fold observée à la résolution 1500 km x 1000 km (cf. figure \ref{...}) : au-delà du faible nombre de blocs disponibles à cette résolution, la difficulté d'un bloc test donné dépend potentiellement de sa position par rapport à des zones de corrélation plus ou moins étendues, et pas seulement de sa distance brute au reste du jeu d'entraînement. 
Une attention particulière sur la dépendance à la lattitude est nécessaire afin d'améliorer le modèle. 


\subsection{Déséquilibre de l'échantillonage}
La figure \ref{label} présente les transects de collecte des données d'acoustique active dans la région d'étude. Les campagnes océanographiques suivent des lignes de transect plutôt qu'une couverture surfacique homogène de la zone d'étude : de larges portions de la zone ne sont jamais échantillonnées, tandis que d'autres sont survolées de façon dense et répétée. Ce déséquilibre a plusieurs conséquences pour l'apprentissage et son évaluation. D'une part, le modèle est entraîné sur une distribution des covariables qui ne représente pas uniformément l'ensemble de la zone d'étude et les régions les plus échantillonnées pèsent alors davantage dans l'ajustement du modèle, au détriment des régions peu visitées, pour lesquelles la prédiction repose sur une extrapolation plus incertaine. D'autre part, ce déséquilibre se répercute directement sur la construction des folds de validation croisée par blocage : le nombre d'observations disponibles par bloc varie fortement selon la résolution et la localisation (de quelques dizaines à plusieurs milliers d'observations par bloc), ce qui explique en partie la variance inter-fold élevée observée aux résolutions les plus strictes, où un unique fold peu représentatif peut disproportionnellement influencer l'estimation moyenne de performance.

Un second aspect du déséquilibre concerne la distribution même de la variable réponse : le NASC présente une distribution asymétrique, avec une majorité de valeurs faibles à modérées et une minorité de valeurs élevées (cf. figure \ref{...}, distribution observée). Les modèles d'arbres, en minimisant une erreur quadratique moyenne, tendent à mieux ajuster les régions denses de la distribution au détriment des valeurs extrêmes, moins représentées, un mécanisme cohérent avec la tendance de sous-prédiction systématique des valeurs élevées observée sur les courbes de calibration (section \ref{...}). Ce déséquilibre de la variable réponse s'ajoute ainsi au déséquilibre spatial de l'échantillonnage, sans qu'il soit possible, dans le cadre de cette étude, de dissocier précisément la part de chacun dans les biais de prédiction observés.

\section{Relation du NASC avec les variables environnementales et phytoplanctoniques}
L'importance des variables environnementales et biologiques dans la structuration du NASC, évalué par RF (voir \ref{label}) est cohérente avec la littérature. Son entre chaque schéma de validation et même chaque modèle \ref{annexe:import_var_all} constitue un signal fort. Les FOD arrivent en première position. Comme discuté en section \ref{annexe:lat_effect}, cela ne signifie pas nécessairement que les conditions thermohalines de l'eau sont les facteurs structurants principaux du NASC, mais peut-être que les FOD capturent juste l'information de latitude. En seconde position se trouve la biomasse totale de phytoplancton, représentée par la chlorophylle-a, ce qui est cohérent avec la littérature : une forte biomasse de phytoplancton est corrélée positivement avec une forte biomasse de NTI \cite{bibid}. Ici, seule l'importance est indiquée et non le sens de la dépendance (positif ou négatif). Les troisième et quatrième variables les plus importantes sont inversées pour 38 et 120 kHz mais ont des ordres de grandeurs proches. Ce sont la fucoxanthine (diatomées) et la Péridinine (Dinoflagellée). Les variables biologiques les plus importantes sont donc celles associées au microphytoplancton. Les nano et picophytoplancton sont moins représentés, traduisant ainsi une moins grande dépendance de la distribution des NTI à ces PFD. Les modèles ne captant pas l'entièreté de la distribution du NASC et prédisant majoritairement les modes principaux \ref{annexe:distrib_train_test}, on peut émettre l'hypothèse que la présence de microphytoplancton (diatomées, dinoflagellées) est un signal très fort pour le modèle : il induit un pic de présence de NTI; l'absence ou la présence de nanophytoplancton influence plus subtilement les distributions et ne sont pas captés par le modèle. 

\section{Incertitudes liées à la nature des données}	
\subsection{Données d'acoustique active}
Plusieurs propriétés des données d'acoustique active sont susceptible de bruiter et biaiser le modèle. Dans un premier temps, il y a le problème de déséquilibre de l'échantillonage présenté en section \ref{label}. D'autre part, les données d'acoustique actives sont bruitées : le phénomène de \textit{bullage} modifie les données à des profondeurs superficielles; des changements de propriétés physique de l'eau (par exemple des changements de densité) peuvent introduire de faux signaux; l'instrumentation électronique voisinant l'échosondeur peut parasiter les données, \textit{etc.}
Par ailleurs, une incertitude non négligeable concerne la métrique de NASC. En effet, le NASC est une intégration sur un ensemble de profondeur. Or, cette profondeur varie selon la fréquence d'échantillonage, les valeurs obtenues ne sont pas comparables entre deux fréquences. D'autre part, les profils de rétro-propagation acoustique présentent des valeurs manquantes pouvant affecter directement la valeur de NASC obtenue. La méthode que nous avons mis en place pour diminuer cet effet (voir section \ref{section:compute_nasc}) présente des faiblesse et mériterait une attention particulière afin de l'améliorer dans de futurs travaux. Ici, le jeu de donnée présentant le plus de valeur manquante et le plus susceptible d'être affecté par ce protocole est la campagne océanographique de 2021 (voir Annexe \ref{annexe:missing_values}).
Enfin, les organismes étudiés par acoustique active présentent des caractéristiques diverses de réflectivité des ondes émises par l'écho-sondeur. Leur composition, orientation dans la colonne d'eau ou bien même la composition en lipides de leurs poche lipidique, variant de manière saisonnière, peuvent affecter le signal reçu \cite{benoitbird2016ecological}. 

\subsection{Données phytoplanctoniques}
Les données phytoplanctoniques issus de l'algorithme PIGMeANN sont des prédictions issues d'un algorithme de deep learning et présentent de ce fait des incertitudes intrinsèques. D'autre part, la couverture nuageuse importante de la région permet au satellite d'échantillonner l'océan seulement 10 à 20\% du temps. Ce phénomène limite fortement le nombre de données disponible et restreint nos modèles d'apprentissage à certaines dates et espaces, ce qui peut induire un certain biais. 
Par ailleurs, l'équivalence entre concentration de pigments et biomasse phytoplanctonique est aussi controversée : les cellules phytoplanctoniques peuvent ajuster leur chlorophylle-a aux ratios de Carbone en réponse au stress environnemental \cite{thomalla2023Widespread}. 

\subsection{Données modèles : FOD et FTLE}
L'approche des FOD pour regrouper des masses d'eau aux propriétés similaires diffère de l'approche usuellement appliquée. En effet, le plus courant est une régionalisation de la zone d'étude selon l'apport d'eau de glace. On mentionne les Coastal and Continental Shelf Zone (CCSZ), Seasonal Ice Zone (SIZ), Permanently Open Ocean Zone (POOZ), Polar Frontal Zone (PFZ) ou Sub-Antarctic Zone (SAZ) \cite{treguer1992}. Plus récemment, les termes de WIZ (Winter Ice Zone) et MIEZ (Maximum Ice Extent Zone) ont également été introduits \cite{adroli2026}. Cette manière de procéder est justifiée par les régimes de productivité très différents selon ces eaux. Travailler avec les FOD permet néanmoins de retrouver cette régionalisation puisque chaque cluster encode les caractéristiques thermohaline de l'eau et donc contient l'information d'adoucissement de l'eau due à la fonte des glaces. La seule exception concerne la SZ : cette zone est marquée par des apports minéraux et en fer d'origine continentale (mise en suspension de sédiments du plateau, fonte glaciaire). Ces processus, propres au domaine côtier, ne sont pas capturés par notre modèle de régionalisation, fondé uniquement sur les propriétés thermohalines des masses d'eau.
D'autre part, les FOD ainsi que les Finite-Time Lyapunov Exponent (FTLE) sont calculés à partir des champs de température, salinité et courants issus de la réanalyse GLORYS12 de Copernicus Marine Service. Cette réanalyse assimile conjointement des données altimétriques, de température de surface satellite, de concentration en glace de mer, ainsi que des profils in situ de température et salinité. Cependant la qualité de cette assimilation dépend directement de la densité d'observations in situ disponibles \cite{lellouche2021copernicus} : dans l'Océan Austral, la couverture en profils Argo et autres plateformes in situ reste nettement plus faible que dans les autres régions. Les champs de température, salinité et courants dans notre zone d'étude sont donc potentiellement moins contraints par l'observation et davantage dépendants de la dynamique propre du modèle NEMO sous-jacent. Ces biais peuvent se traduire par exemple par des décalages dans la délimitation spatiale des domaines fonctionnels identifiés. Enfin, la réanalyse fournit des champs quotidiens, ce qui limite la capacité du modèle à capturer la variabilité sub-journalière susceptible d'affecter transitoirement la structuration des masses d'eau et la dynamique de mélange de la région étudiée.




\section{Perspectives}
Une approche intégrant la phénologie du phytoplancton et les délais (connus et documentés \cite{nicholson2025satellite}, \cite{kauko2021phenology}, \cite{lampe2026differences}) de la colonisation des NTI suite à des blooms rapprocherait le modèle de la réalité biologique et permettrait potentiellement d'améliorer les prédictions. Il serait intéressant d'intégrer cette approche dans le cadre de futurs travaux. Une attention particulière sur le respect de la stationnarité permettrait d'évaluer avec d'avantage de précision l'importance des variables mais également d'adapter plus finement les résolutions de validation croisée et donc d'améliorer les performances des modèles. Elargir l'ensemble de données et sa couverture spatiale pemettrait d'affiner les prédictions et de les rendre plus fiable.  

